% ============================================================
%  BFPME_2 — RAPPORT DE SYNTHÈSE (PHASE 2 : DONNÉES RÉELLES)
%  Chapitre 1 : jeux de données externes (réels, publics)
%  Chapitre 2 : données réelles de la BFPME
%  Tous les chiffres proviennent des résultats calculés du projet bfpme_2.
% ============================================================
\documentclass[a4paper,11pt]{report}

% ── Encodage et langue ───────────────────────────────────────
\usepackage[utf8]{inputenc}
\usepackage[T1]{fontenc}
\usepackage[french]{babel}
\usepackage{lmodern}
\usepackage{microtype}

% ── Mise en page ─────────────────────────────────────────────
\usepackage[top=2.5cm, bottom=2.5cm, left=2.5cm, right=2.5cm]{geometry}
\usepackage{setspace}
\onehalfspacing

% ── Couleurs ─────────────────────────────────────────────────
\usepackage[dvipsnames,svgnames,table]{xcolor}
% Tout le texte en noir : BfRed et DarkGray redéfinis en noir pur.
% (Les gris clairs ne servent qu'au remplissage des figures.)
\definecolor{BfRed}{RGB}{0, 0, 0}
\definecolor{DarkGray}{RGB}{0, 0, 0}
\definecolor{LightGray}{RGB}{235, 235, 235}
\definecolor{MidGray}{RGB}{150, 150, 150}
\definecolor{GoodGreen}{RGB}{0, 0, 0}

% ── Graphiques ───────────────────────────────────────────────
\usepackage{float}
\usepackage{tikz}
\usetikzlibrary{arrows.meta, positioning, shapes.geometric, calc, fit}

% ── Titres ───────────────────────────────────────────────────
\usepackage{titlesec}
\titleformat{\chapter}[display]
  {\normalfont\bfseries\color{BfRed}}
  {\Large CHAPITRE \thechapter}{6pt}{\Huge}
\titlespacing*{\chapter}{0pt}{-20pt}{24pt}
\titleformat{\section}{\normalfont\Large\bfseries\color{BfRed}}{\thesection}{0.6em}{}
\titleformat{\subsection}{\normalfont\large\bfseries\color{DarkGray}}{\thesubsection}{0.6em}{}

% ── Tableaux et listes ───────────────────────────────────────
\usepackage{booktabs}
\usepackage{tabularx}
\usepackage{array}
\usepackage{multirow}
\usepackage{enumitem}
\setlist[itemize]{leftmargin=1.6em, itemsep=2pt, topsep=3pt}
\newcolumntype{Y}{>{\raggedright\arraybackslash}X}
\newcolumntype{C}{>{\centering\arraybackslash}X}

% ── Mathématiques ────────────────────────────────────────────
\usepackage{amsmath, amssymb}

% ── En-tête / pied ───────────────────────────────────────────
\usepackage{fancyhdr}
\pagestyle{fancy}\fancyhf{}
\fancyhead[L]{\small\textit{BFPME\_2 — Synthèse phase~2}}
\fancyhead[R]{\small Données réelles \& externes}
\fancyfoot[C]{\thepage}
\renewcommand{\headrulewidth}{0.4pt}
\setlength{\headheight}{14.5pt}

% ── Légendes ─────────────────────────────────────────────────
\usepackage{caption}
\captionsetup{labelsep=endash, font=small, labelfont=bf, justification=centering}

\usepackage{hyperref}
\hypersetup{colorlinks=true, linkcolor=BfRed, urlcolor=BfRed, citecolor=BfRed}

% ── Barre de métrique horizontale : \aucbar{y}{label}{val}{couleur} ──
\newcommand{\aucbar}[4]{%
  \node[anchor=east, font=\small] at (0,#1) {#2};
  \draw[fill=LightGray, draw=MidGray] (0.2,{#1-0.2}) rectangle (6.2,{#1+0.2});
  \draw[fill=#4, draw=#4] (0.2,{#1-0.2}) rectangle ({0.2+6.0*#3},{#1+0.2});
  \node[anchor=west, font=\small\bfseries] at (6.35,#1) {#3};
}

\begin{document}

% ── Page de titre ────────────────────────────────────────────
\begin{titlepage}
\centering
\vspace*{2cm}
{\Huge\bfseries\color{BfRed} BFPME\_2\par}
\vspace{0.4cm}
{\Large Confrontation aux données réelles~: jeux externes et dossiers BFPME\par}
\vspace{1.2cm}
\rule{0.8\textwidth}{1.5pt}\par
\vspace{0.8cm}
{\LARGE\bfseries Rapport de synthèse — Phase~2\par}
\vspace{0.3cm}
{\large Chapitre~1 : jeux de données externes\\Chapitre~2 : données réelles de la BFPME\par}
\vspace{0.6cm}
{\large\itshape Travaux, méthodes, résultats et illustrations\par}
\vspace{1cm}
\rule{0.8\textwidth}{1.5pt}\par
\vfill
{\large Société d'accueil~: MPSoft\par}
{\large Projet de fin d'études\par}
\vspace{1cm}
\end{titlepage}

\tableofcontents
\newpage

% ── Introduction ─────────────────────────────────────────────
\chapter*{Introduction}
\addcontentsline{toc}{chapter}{Introduction}

La phase~1 (rapport \textit{BFPME}) a validé, sur données \emph{synthétiques}, une chaîne
complète de prédiction, d'explication et de test de stress du risque de crédit PME. La phase~2
confronte cette approche au \textbf{réel}, selon deux axes traités chacun dans un chapitre~:

\begin{itemize}
  \item \textbf{Chapitre~1 — Jeux de données externes.} La méthodologie BFPME est répliquée à
        l'identique sur deux jeux de crédit \emph{publics et réels} (Taïwan, Home Credit) afin
        d'éprouver sa transférabilité hors du cadre du projet.
  \item \textbf{Chapitre~2 — Données réelles de la BFPME.} Les 1\,226~dossiers réels de la banque
        servent d'abord de \emph{jeu de validation externe} pour mesurer si le modèle synthétique
        généralise, puis de base d'un modèle entraîné \emph{directement} sur le réel.
\end{itemize}

\noindent Chaque chapitre suit la même trame~: objectifs, méthodes, résultats et illustrations.
Le fil conducteur est la \textbf{crédibilité}~: un score élevé n'a de valeur que s'il survit au
contact des dossiers réels.

\begin{figure}[H]
\centering
\begin{tikzpicture}[node distance=8mm and 12mm,
   src/.style={draw=DarkGray, thick, rounded corners=2pt, fill=LightGray,
               text width=3.2cm, minimum height=1.1cm, align=center, font=\scriptsize},
   core/.style={draw=BfRed, very thick, rounded corners=5pt, fill=white,
                text width=3cm, minimum height=1.3cm, align=center, font=\scriptsize\bfseries},
   fl/.style={-{Stealth[length=2.2mm]}, thick, DarkGray}]
  \node[src] (ext) {Jeux externes réels\\Taïwan · Home Credit\\(chap.~1)};
  \node[src, below=of ext] (real) {Dossiers réels BFPME\\1\,226 lignes (chap.~2)};
  \node[core, right=of $(ext)!0.5!(real)$] (pipe) {Méthodologie BFPME\\(prétraitement, modèle, XAI)};
  \node[src, right=of pipe] (val) {Validation \&\\gouvernance};
  \draw[fl] (ext) -- (pipe);
  \draw[fl] (real) -- (pipe);
  \draw[fl] (pipe) -- (val);
\end{tikzpicture}
\caption{Deux sources de données réelles éprouvant la même méthodologie}
\end{figure}

% ============================================================
\chapter{Jeux de données externes}
% ============================================================

\section{Objectifs}

Avant d'exploiter les données internes de la banque, la méthodologie BFPME (préparation,
\textit{feature engineering}, modélisation XGBoost / RandomForest, comparaison par bootstrap
apparié) est \textbf{répliquée à l'identique} sur deux jeux de crédit publics et réels. Objectif~:
vérifier que le pipeline se transpose à d'autres populations et produit des performances
crédibles sur des données non issues du générateur BFPME.

\section{Méthodes}

Les deux jeux suivent exactement la chaîne des notebooks BFPME (\texttt{01\_préparation},
\texttt{02\_xgboost}, \texttt{03\_randomforest}, \texttt{04\_comparaison})~: nettoyage,
ingénierie de variables métier, préprocessing sans fuite de cible, évaluation en validation
croisée \textit{out-of-fold} (AUC / Gini / KS / PR-AUC), hold-out avec intervalle de confiance
bootstrap, puis désignation du modèle retenu.

\begin{table}[H]
\centering
\caption{Jeux de données externes réels étudiés}
\renewcommand{\arraystretch}{1.3}
\begin{tabularx}{\textwidth}{@{}l Y C C@{}}
\toprule
\textbf{Jeu} & \textbf{Nature} & \textbf{\# clients} & \textbf{Taux de défaut} \\
\midrule
Credit Card Taïwan (UCI) & Table unique, 6 mois d'historique de paiement & 30\,000 & $\approx$ 22\,\% \\
Home Credit (Kaggle)     & 1 table principale + 6 tables satellites (relationnel) & 307\,511 & $\approx$ 8\,\% \\
\bottomrule
\end{tabularx}
\end{table}

\section{Résultats}

Sur les deux jeux, XGBoost l'emporte sur RandomForest (bootstrap apparié, écart d'AUC
significatif), avec des AUC de hold-out de l'ordre de \textbf{0{,}78}
(tableau~\ref{tab:externes}, figure~\ref{fig:externes})~— des niveaux réalistes et cohérents
avec l'état de l'art du \textit{credit scoring}.

\begin{table}[H]
\centering
\caption{Résultats des jeux externes (hold-out ; modèle retenu en gras)}
\label{tab:externes}
\renewcommand{\arraystretch}{1.3}
\begin{tabularx}{\textwidth}{@{}l l C C C@{}}
\toprule
\textbf{Jeu} & \textbf{Modèle} & \textbf{AUC} & \textbf{Gini} & \textbf{KS} \\
\midrule
\multirow{2}{*}{Credit Card Taïwan}
  & \textbf{XGBoost}      & \textbf{0{,}782} & 0{,}565 & 0{,}431 \\
  & RandomForest          & 0{,}779 & 0{,}557 & 0{,}428 \\
\midrule
\multirow{2}{*}{Home Credit}
  & \textbf{XGBoost}      & \textbf{0{,}787} & 0{,}575 & 0{,}439 \\
  & RandomForest          & 0{,}768 & 0{,}536 & 0{,}403 \\
\bottomrule
\end{tabularx}
\end{table}

\begin{figure}[H]
\centering
\begin{tikzpicture}
  \aucbar{0}{Taïwan · RandomForest}{0.779}{MidGray}
  \aucbar{-0.7}{Taïwan · XGBoost}{0.782}{BfRed}
  \aucbar{-1.6}{Home Credit · RandomForest}{0.768}{MidGray}
  \aucbar{-2.3}{Home Credit · XGBoost}{0.787}{BfRed}
\end{tikzpicture}
\caption{AUC de hold-out sur les jeux externes réels (XGBoost en rouge)}
\label{fig:externes}
\end{figure}

\noindent\textbf{Enseignement.} La méthodologie BFPME, appliquée sans modification à des données
réelles publiques, atteint des performances solides et stables. Le pipeline est donc
techniquement transférable~; reste à savoir s'il se transfère à la population \emph{spécifique}
de la BFPME — c'est l'objet du chapitre~2.

% ============================================================
\chapter{Données réelles de la BFPME}
% ============================================================

\section{Objectifs}

Ce chapitre exploite les \textbf{1\,226~dossiers réels} de la BFPME (fichier \texttt{Data.xlsx},
taux de défaut réel \textbf{65,0\,\%}). Il répond à deux questions distinctes~:
\begin{enumerate}[leftmargin=1.8em]
  \item Le modèle entraîné sur les données \emph{synthétiques} de la phase~1 généralise-t-il aux
        dossiers \emph{réels}~? (validation externe)
  \item Que vaut un modèle entraîné \emph{directement} sur les données réelles~?
\end{enumerate}

\section{Validation externe du modèle synthétique}

\subsection{Méthode}

Plutôt que de \emph{fusionner} données synthétiques et réelles — impossible faute de clé
commune et de schéma partagé — les dossiers réels sont utilisés comme \textbf{jeu de test
indépendant} (validation externe). On ne conserve que les variables présentes des deux côtés et
sans fuite de cible~: \texttt{log\_montant\_pret}, \texttt{log\_cout\_total},
\texttt{ratio\_financement\_externe}, \texttt{secteur}, \texttt{objectif}. Trois mesures sont
comparées, pour deux familles de modèles (régression logistique et gradient boosting)~:

\begin{itemize}
  \item \textbf{A} — AUC interne sur données synthétiques (le score « qu'on rapporterait »)~;
  \item \textbf{B} — AUC sur les dossiers réels (généralisation effective)~;
  \item \textbf{C} — AUC d'un modèle entraîné \emph{uniquement} sur le réel, sur ces mêmes
        5~variables (plafond atteignable avec ce périmètre).
\end{itemize}

\subsection{Résultats : l'effondrement sur données réelles}

\begin{table}[H]
\centering
\caption{Validation externe — synthétique vs.\ réel (5 variables partagées)}
\label{tab:benchmark}
\renewcommand{\arraystretch}{1.3}
\begin{tabularx}{\textwidth}{@{}Y C C@{}}
\toprule
\textbf{Mesure} & \textbf{Régression logistique} & \textbf{Gradient Boosting} \\
\midrule
A. AUC interne (synthétique) & \textbf{0{,}877} & \textbf{0{,}879} \\
B. AUC externe (réel)        & 0{,}527 & 0{,}507 \\
C. AUC réel seul (5 var.)    & 0{,}522 & 0{,}545 \\
\midrule
Écart de réalité (A $-$ B)   & 0{,}350 & 0{,}371 \\
KS sur le réel (B)           & 0{,}088 & 0{,}058 \\
\bottomrule
\end{tabularx}
\end{table}

Le constat est \textbf{identique pour les deux familles de modèles} — c'est donc une propriété
des \emph{données}, non d'un algorithme~:
\begin{itemize}
  \item sur données synthétiques, le modèle paraît excellent (AUC $\approx$ 0{,}88)~;
  \item sur données réelles, il s'effondre à AUC $\approx$ 0{,}51 — l'équivalent d'un tirage à
        pile ou face. Le modèle a appris les \emph{règles du générateur}, pas le risque de crédit~;
  \item le plafond réel avec ces 5~variables (C $\approx$ 0{,}52) est lui aussi proche du hasard~:
        ces variables partagées ne sont, à elles seules, \textbf{pas prédictives} du défaut réel.
\end{itemize}

\begin{figure}[H]
\centering
\begin{tikzpicture}
  \aucbar{0}{A · synthétique (interne)}{0.877}{BfRed}
  \aucbar{-0.7}{B · réel (validation externe)}{0.527}{MidGray}
  \aucbar{-1.4}{C · réel seul (5 var.)}{0.522}{MidGray}
  \draw[dashed, DarkGray] ({0.2+6.0*0.5},0.3) -- ({0.2+6.0*0.5},-1.7)
    node[below, font=\scriptsize\itshape]{hasard (0{,}50)};
\end{tikzpicture}
\caption{Écart de réalité (régression logistique)~: le score chute de 0{,}88 à 0{,}53}
\label{fig:threeworlds}
\end{figure}

Le tri par PD prédite le rend concret~: le décile le \emph{moins} risqué (PD moyenne 6\,\%)
fait tout de même défaut à \textbf{69,9\,\%} — autant que le décile le plus risqué. Le modèle
synthétique ne \emph{classe} donc pas les emprunteurs réels.

\section{Modèle entraîné directement sur le réel}

\subsection{Méthode}

Le second volet entraîne un modèle \emph{directement} sur les 1\,226~dossiers réels
(\texttt{dataset\_clean.csv}), cette fois avec l'ensemble des \textbf{40~variables} réelles
disponibles — et non plus les seules 5~variables partagées. Protocole~: validation croisée pour
la sélection, hold-out final (980~dossiers d'apprentissage, 246~de test), et exploration de
stratégies d'usage du signal de score BFPME.

\subsection{Résultats}

Avec ses 40~variables natives, le modèle réel dépasse nettement le plafond des 5~variables
partagées (AUC 0{,}52) et atteint des niveaux exploitables (tableau~\ref{tab:realmodels}). Le
RandomForest réglé est retenu~; sur le hold-out final il obtient
\textbf{AUC~=~0{,}750}, \textbf{KS~=~0{,}430}, rappel~=~0{,}906.

\begin{table}[H]
\centering
\caption{Comparaison des modèles sur données réelles (40 variables, validation croisée)}
\label{tab:realmodels}
\renewcommand{\arraystretch}{1.3}
\begin{tabularx}{\textwidth}{@{}l C C C C@{}}
\toprule
\textbf{Modèle} & \textbf{AUC} & \textbf{Gini} & \textbf{KS} & \textbf{F1} \\
\midrule
\textbf{RandomForest réglé} & \textbf{0{,}719} & 0{,}438 & 0{,}317 & 0{,}782 \\
RandomForest                & 0{,}716 & 0{,}431 & 0{,}325 & 0{,}796 \\
XGBoost                     & 0{,}702 & 0{,}404 & 0{,}289 & 0{,}751 \\
Régression logistique (base)& 0{,}675 & 0{,}351 & 0{,}275 & 0{,}720 \\
\bottomrule
\end{tabularx}
\end{table}

\begin{table}[H]
\centering
\caption{Modèle retenu sur le hold-out réel (RandomForest, \texttt{model\_metrics.json})}
\renewcommand{\arraystretch}{1.3}
\begin{tabularx}{\textwidth}{@{}Y C C C C C@{}}
\toprule
\textbf{AUC} & \textbf{Gini} & \textbf{KS} & \textbf{Exactitude} & \textbf{Précision} & \textbf{Rappel} \\
\midrule
0{,}750 & 0{,}501 & 0{,}430 & 0{,}768 & 0{,}775 & 0{,}906 \\
\bottomrule
\end{tabularx}
\end{table}

\noindent Les stratégies d'intégration du score BFPME officiel confirment un léger apport~: à
partir des seules données (AUC 0{,}721), l'ajout de l'indicateur de présence de note porte l'AUC
à 0{,}723. Le signal reste modeste, ce qui rejoint le diagnostic de la validation externe~: le
pouvoir prédictif le plus fort réside dans des variables \emph{absentes} du fichier réel actuel
(sous-facteurs \texttt{SF\_1..SF\_62}, ratios financiers granulaires, historique comportemental).

% ── Conclusion ───────────────────────────────────────────────
\chapter*{Conclusion}
\addcontentsline{toc}{chapter}{Conclusion}

La phase~2 apporte la preuve, chiffrée, de deux faits complémentaires. D'une part, la
méthodologie BFPME est \textbf{techniquement solide}~: appliquée telle quelle à des jeux de
crédit réels publics (Taïwan, Home Credit), elle atteint des AUC de l'ordre de 0{,}78. D'autre
part, le modèle de la phase~1 \textbf{ne généralise pas} aux dossiers réels de la BFPME~: son
AUC chute de 0{,}88 à 0{,}51, révélant que les performances synthétiques mesuraient la cohérence
du générateur, non le risque de crédit réel.

Confronté directement aux 1\,226~dossiers réels avec leurs 40~variables, un modèle atteint en
revanche une AUC de \textbf{0{,}72–0{,}75}, exploitable mais bornée par la richesse limitée du
fichier actuel. La conclusion opérationnelle est nette~: la valeur des données synthétiques était
de valider le \emph{pipeline}~; la valeur des données réelles est de fournir la \emph{seule} base
fiable pour un modèle déployable. La suite consiste à récupérer, dans le système d'information de
la banque, les variables réellement prédictives (sous-facteurs de notation, ratios financiers
détaillés, historique de paiement) absentes du fichier actuel, et à ré-entraîner le modèle sur
cette base enrichie avant tout déploiement.

\end{document}
